\scanpage{8}
\mspage{\textcircled{3}}
\[
P\otimes_kQ\hookrightarrow(P\otimes_kQ)^{\unsure{\&}}
\simeq\underline{\operatorname{Bil}}_k(P^\circ,Q^\circ;(Ab_k))
\]
est donné par
\[
(*)\qquad a\otimes_kb\longmapsto
\big((x,y)\longmapsto\mathrm{Hom}_P(x,a)\otimes_k\mathrm{Hom}_Q(y,b)\big)
\]
ce qui donne une construction \unsure{directe} de
$P\otimes_kQ$ comme la sous-cat. pleine de
$\underline{\operatorname{Bil}}_k(P^\circ,Q^\circ;Ab_k)$
formée des sommes finies des foncteurs de la forme $(*)$
(\uline{ou} la cat. str. pleine correspondante).

\margin{\figure{\begin{tikzpicture}\draw[->] (0,0)--(1,.5);\end{tikzpicture}}}
\uline{Application 1} Soit $M\simeq Q^{\unsure{\&}}$ un topos $k$-abélien \unsure{alg.},
$M^\vee\simeq Q^{\circ\unsure{\&}}$ ``le'' topos $k$-abélien \unsure{alg.} dual
\margin{(de l'existence de $P\otimes_kQ$, \deleted{\unsure{et}} et du fait qu'il est petit si $P,Q$ le sont)}
Considérons $P_M^{\unsure{\&}}$ et $P_{M^\vee}^{\circ\unsure{\&}}$
($\simeq\mathcal P_M$ et $(\mathcal P^\vee)_{M^\vee}$),
ce sont des topos $k$-abéliens \unsure{alg.}, et en dualité.

En effet,
\[
\begin{aligned}
P_M^{\unsure{\&}}&=\underline{\mathrm{Hom}}_k(P^\circ,\underline{\mathrm{Hom}}_k(Q^\circ,Ab_k))\\
&\simeq\underline{\mathrm{Hom}}_k(\underbrace{P^\circ\otimes_kQ^\circ}_{\rotatebox{90}{$\simeq$}},Ab_k)\\
&\simeq\underline{\mathrm{Hom}}_k(\overbrace{(P\otimes_kQ)^\circ},Ab_k)
=(P\otimes_kQ)^{\unsure{\&}}.
\end{aligned}
\]
donc topos $k$-abélien \unsure{alg.}, de $=$
\[
\begin{aligned}
P_{M^\vee}^{\circ\unsure{\&}}&\simeq\underline{\mathrm{Hom}}_k(
\underbrace{(P^\circ\otimes_kQ^\circ)^\circ},Ab_k)\\
&\simeq\underline{\mathrm{Hom}}(P\otimes_kQ,Ab_k)
=(P\otimes_kQ)^{\circ\unsure{\&}}
\end{aligned}
\]
\uline{Appl. 2}
\[
\underline{\operatorname{Bil}}_{k!!}(P^{\unsure{\&}},Q^{\unsure{\&}};M)
\xleftarrow{\simeq}\deleted{\unsure{Bil}}
\underline{\mathrm{Hom}}_{k!}((P\otimes_kQ)^{\unsure{\&}},M)
\quad(\simeq\operatorname{Bil}_k(P,Q;\unsure{M}))
\]
\margin{si $M$ $k$-add. avec $\varinjlim$}
\margin{$\unsure{*}$}
i.e. $(P^{\deleted{\unsure{\&}}}\otimes_kQ)^{\unsure{\&}}$
joue le rôle de prod. tens. des topos $k$-abél. $P^{\unsure{\&}},Q^{\unsure{\&}}$,
parmi les cat. $k$-additives st.\unsure{bles} par $\varinjlim$
% archive foliation: 7
